\section[C++ network experiments]{C++ network experiments: assignment probabilities, graph identity and distributed state}
\label{sec:exec25-cxx-network}

A bank evaluating interventions across connected borrowers needs more than a fast graph engine. Its experiment must distinguish the assignment law, the exposure event that makes an observed outcome informative about a policy, and the numerical representation used to compute that event. A program may preserve adjacency while changing weights, reuse a random stream across workers, or return a number when one exposure arm is empty. Each change can alter the scientific result. This comparison examines the RGCR research implementation alongside NetworKit, SNAP, Repast HPC and the C-native igraph library as an interoperability alternative. The first supplies a specific experimental-design implementation; the others supply components whose scientific role must be declared.~\citep{exec25_cnet_rgcr,exec25_cnet_networkit,exec25_cnet_snap,exec25_cnet_repast,exec25_cnet_igraph}

The comparison binds 64 selected software, license and configuration originals at revisions before the 7 October 2026 cutoff. Six bounded repository queries produce 52 distinct leads, with unreturned search tails and uninspected candidates. They cannot establish complete C++ coverage or absence of other causal software. NetworKit's selected metadata says 11.2.1, but its source build sets C++20 and CMake 3.14; those exact requirements take precedence over a less demanding README snippet for this revision. RGCR's build instead references C++11 and a user-specific SNAP 4.1 installation. The independently inspected SNAP revision is not evidence that this older dependency has been reproduced.~\citep{exec25_cnet_choice,exec25_cnet_networkit,exec25_cnet_rgcr}

The methodological discussion uses the genuine 59-page arXiv version of Ugander and Yin, submitted on 4 September 2020, whose title page bears 7 September 2020. The later publisher edition was not inspected. Version identity matters: the argument here cites the selected author text and distinguishes its statements from the project's finite calculations and the implementation's conventions. No acquired program, simulator, native library or upstream test was executed.~\citep{exec25_cnet_paper,exec25_cnet_choice}

\subsection{Randomized clustering changes both design and analysis}

Let $G=(V,E)$ be a fixed undirected, unweighted interference graph with $n=|V|$ borrowers. For binary assignment vector $Z$, let $Y_i(Z)$ be borrower $i$'s potential outcome and define the global average treatment effect
\[
 \tau=\frac{1}{n}\sum_{i=1}^n\{Y_i(\mathbf{1})-Y_i(\mathbf{0})\}.
\]
The full-neighborhood exposure event $E_i^a$ occurs when borrower $i$ and every neighbor receive assignment $a\in\{0,1\}$. Using the observed outcome on that event requires the substantive assumption that it equals $Y_i(\mathbf{a})$. An adjacency file cannot establish this exposure model. A graph of transactions, for example, need not include all shared employers, household obligations or unobserved channels relevant to a servicing response.~\citep{exec25_cnet_paper,exec25_cnet_choice}

In randomized graph cluster randomization, a partition $C$ is sampled from a declared distribution $P$, followed by treatment assignment at the cluster level. The analysis uses the marginal exposure probability
\[
 \pi_i^a=\Pr(E_i^a\mid P)=\sum_c P(c)\Pr(E_i^a\mid C=c),
 \qquad
 \widehat\tau_{\mathrm{HT}}=
 \frac{1}{n}\sum_i\left\{
 \frac{\mathbf{1}\{E_i^1\}Y_i(\mathbf{1})}{\pi_i^1}
 -\frac{\mathbf{1}\{E_i^0\}Y_i(\mathbf{0})}{\pi_i^0}
 \right\}.
\]
The selected paper derives this marginalization explicitly. Conditioning only on the realized partition gives a different weighting scheme. Likewise, averaging inverse conditional probabilities does not compute the inverse marginal probability. If two equally likely partitions give exposure probabilities $1/4$ and $1/2$, the marginal probability is $3/8$ and its inverse is $8/3$, whereas the average of the two inverse probabilities is three. The adapter should bind the partition generator and its distribution alongside the realized partition.~\citep{exec25_cnet_paper,exec25_cnet_choice}

The source implements several partition rules, including neighborhood maxima and randomized nets. For independent cluster assignment it constructs each borrower's set of adjacent clusters, including its own, and evaluates exposure probability as $2^{-k_i}$, where $k_i$ is that set's size. It also accumulates same-arm and opposite-arm joint exposure probabilities. This is substantive support for interference-aware design analysis. It is specifically tied to the inspected half-probability assignment: the treatment simulator uses a Bernoulli distribution with parameter $1/2$. Changing only a user-facing field to $p=1/5$ would not change these expressions. For a neighborhood touching two independently assigned clusters, treatment and control exposure probabilities would be $1/25$ and $16/25$, respectively, rather than $1/4$ for each.~\citep{exec25_cnet_rgcr,exec25_cnet_choice}

The implementation's complete-randomization mode is also specific. It orders clusters by decreasing size, pairs them, and gives each pair opposite assignments; an unpaired final cluster receives its own assignment. If a borrower's neighborhood intersects both members of a pair, full exposure is impossible under that partition. The helper returns zero for this conflict, and the analysis records excluded target contributions. This does not restore positivity. Equal-size clusters have no explicit secondary key in the selected comparator, so the realized pairing and its provenance should be saved instead of reconstructed from sizes alone. The declared law is not arbitrary complete randomization or a general bank budget-constrained design.~\citep{exec25_cnet_rgcr,exec25_cnet_choice}

The paper offers a route to controlling rare exposure under its selected random partitions. For its randomized 3-net or 1-hop-max constructions, the reported treatment-exposure lower bound is $p/|B_2(i)|$, where $B_2(i)$ is the set of vertices within two hops, including $i$. Its Monte Carlo probability estimator averages conditional probabilities over sampled partitions; the stated relative mean-square error bound is $|B_2(i)|/(Kp)$ for $K$ samples. The paper also describes stratifying the random rankings so every node is favored, preventing zero estimated exposure in that construction. These statements concern the specified graph and randomization model. They do not validate an unknown bank interference graph or make reciprocal estimated probabilities unbiased. Probability approximation and its contribution to downstream error remain explicit qualification tasks.~\citep{exec25_cnet_paper,exec25_cnet_choice}

For 1-hop-max and bounded nonnegative responses, the selected text further bounds mean-outcome HT variance using $|B_4(i)|$ and $\pi_i^1$; under a restricted-growth condition it obtains a bound proportional to $n^{-1}(d_{\max}+1)^2\kappa^4/p$. Here $d_{\max}$ is maximum degree and $\kappa$ controls neighborhood growth. A bank cannot claim decreasing uncertainty merely by increasing its row count: degree, growth, exposure probabilities and outcome bounds matter. The text expressly distinguishes the available 1-hop-max variance proof from the randomized 3-net case. No unrestricted network-uncertainty guarantee is imported.~\citep{exec25_cnet_paper}

\subsection{An empty realized exposure arm is different from positive design support}

Positive $\pi_i^a$ for every borrower does not guarantee that both exposure arms contain observed borrowers in each experiment. This distinction matters for the ratio form of the H\'ajek estimator,
\[
 \widehat\mu_a^{\mathrm{H}}=
 \frac{\sum_i\mathbf{1}\{E_i^a\}Y_i(\mathbf a)/\pi_i^a}
      {\sum_i\mathbf{1}\{E_i^a\}/\pi_i^a}.
\]
The selected implementation initializes each arm mean to zero and updates it only when its denominator is positive. This defines a numerical extension for an empty arm. It is not the same expression as an everywhere-defined ratio, and zero is not an observed control outcome. The proposed result type should retain both exposed counts, both denominators and the convention or refusal applied.~\citep{exec25_cnet_rgcr,exec25_cnet_choice}

Consider three borrowers in one cluster, with potential outcomes $Y(\mathbf0)=(1,2,3)$ and $Y(\mathbf1)=(2,3,4)$ in a hypothetical monetary unit. The true average effect is one. Assigning the whole cluster with probability $1/2$ gives HT effect estimates $-4$ and $6$, each with probability $1/2$, so their mean is one. The source-style zero extension of H\'ajek gives $-2$ when everyone is controlled and $3$ when everyone is treated, yielding mean $1/2$. Every node has positive exposure probability in both arms, yet one realized arm is always empty. This exact example concerns the stated convention, not a native software run.~\citep{exec25_cnet_choice}

The same issue occurs in a small nontrivial randomized partition example. On the path $1-2-3$, enumerate the six strict orderings of independent continuous node scores and assign each node to its largest one-hop score. Across the induced partitions and independent half-probability cluster assignments, there are 20 weighted cases. The marginal full-exposure probabilities in each arm are $(5/12,1/3,5/12)$. With the preceding potential outcomes, exact enumeration gives HT mean one and variance $2242/75$; an independent joint-probability covariance expansion gives the same variance. The zero-extended H\'ajek mean is again $1/2$. These calculations both verify the intended exposure law and expose the effect of the empty-arm convention.~\citep{exec25_cnet_choice}

The selected 2020 text states a constant-effect unbiasedness result for H\'ajek and proves it by separating a constant treatment effect from each realized ratio. That algebra requires a nonzero realized denominator. The proof does not establish the source code's zero-extension guarantee on the assignments just enumerated. The proposal therefore retains this as a precise qualification boundary. Possible routes include a typed unavailable result, an explicitly redefined estimand or estimator, or an assignment design that guarantees the required denominators. Conditioning on successful experiments would itself change the sampling law and cannot be introduced without analysis. The later publisher version remains a separate source-review obligation.~\citep{exec25_cnet_paper,exec25_cnet_choice}

\subsection{RGCR integration requires more than recompilation}

The RGCR loader first looks for a binary graph cache and otherwise reads an edge list and writes that cache. It removes self-edges and requires consecutive node identifiers, connectivity and a loopless graph. These are explicit, useful restrictions, but they exclude some banking populations unless an adapter makes a justified transformation. A disconnected borrower network should not be made connected by inventing edges. A cached graph should not silently survive a changed source edge list. Stable borrower-to-node mappings, graph hashes and transformation receipts belong in the request.~\citep{exec25_cnet_rgcr,exec25_cnet_choice}

Randomness and input validation are additional gates. The inspected utility seeds the C generator from the clock; clustering and assignment use draws from that generator, while parallel simulation loops create work across threads. A safe analysis distinguishes reproducible seed input, draw allocation, thread scheduling and statistical independence. The build disables assertions, so an assertion in a parser or weight check is not a production rejection path. The weighted-score helper evaluates a logarithm of a draw that can be zero and divides by supplied weights; positive finite weights, endpoint handling and finite-score ties require explicit tests. A mixing-count helper also accesses the last element of a vector that can be empty for requests below ten. A proposed wrapper should reject unsupported counts before entering that path.~\citep{exec25_cnet_rgcr}

The simulation compares complementary assignments using the same partition. That pairing can reduce or change variability, but it means the two outputs are not automatically independent observations. The reported standard-error arithmetic needs a comparison against replicate-level or pair-level uncertainty on frozen small problems. Analytic variance routines also use supplied potential responses and joint exposure probabilities: this is valuable for known-law public evaluation, not an immediately available variance oracle for observed bank outcomes.~\citep{exec25_cnet_rgcr,exec25_cnet_choice}

The proposed decision is conditional algorithmic reference and controlled integration after these contracts are tested. No license-named file appears in the complete selected RGCR tree, so public availability is not treated as a redistribution grant. The old SNAP dependency, full artifacts and rights must be resolved before adoption. An independent reimplementation would need its own rights review and semantic tests; changing language does not remove the assignment, exposure or statistical obligations.~\citep{exec25_cnet_choice}

\subsection{NetworKit and SNAP: preserve identities, multiplicity and weights}

NetworKit's selected graph API supports optional weights and directed graphs but warns that multiple edges are not handled consistently. Its duplicate check is optional and disabled by default. The implementation confirms that enabling the check rejects a repeated edge; otherwise it appends adjacency entries. A borrower graph derived from repeated transactions must therefore choose an aggregation rule before admission. Counting transactions, summing monetary weights and deduplicating relationships define different exposure operators. The adapter should never rely on an unsupported duplicate representation merely because insertion succeeds.~\citep{exec25_cnet_networkit}

Deleting a NetworKit node leaves a hole in its identifier space; the actual node count and the upper identifier bound are separate values. SNAP likewise maintains a maximum assigned identifier bound independently of the number of remaining nodes. In an authored example with surviving IDs zero and two, there are two units, not three. If each has outcome one, dividing a zero-filled three-slot vector by its length produces $2/3$ instead of one. Compacting is a valid representation change only when outcomes, assignments, graph structure and returned results all use the saved forward and inverse maps.~\citep{exec25_cnet_networkit,exec25_cnet_snap,exec25_cnet_choice}

NetworKit's selected induced-subgraph and remapping constructions carry edge weights explicitly. Its restoration helper, however, traverses neighbors and calls edge insertion without the weight argument, whose default is one. Thus a node-ID restoration is not a guaranteed weighted-graph round trip in this selected path. If two incoming relationships have weights nine and one and neighbor treatments one and zero, the normalized weighted exposure is $9/10$; replacing both weights by one changes it to $1/2$. A public round-trip test should check node IDs, edge direction, multiplicity, weights and attributes separately. The finding is tied to this helper and revision, not asserted for every graph operation.~\citep{exec25_cnet_networkit,exec25_cnet_choice}

SNAP offers simple undirected and directed graph types and a separate directed multigraph. In the selected simple graph, duplicate insertion returns $-2$, while a successful insertion without an edge ID returns $-1$. Treating every negative return as failure would therefore misclassify successful operations. Its breadth-first traversal takes explicit incoming and outgoing direction flags and increments hops. The proposed wrapper should translate these API-specific conventions into typed outcomes and a declared distance definition. A weighted credit-exposure distance cannot be inferred from a hop traversal.~\citep{exec25_cnet_snap}

The libraries also expose different random-state contracts. SNAP's selected random generator treats zero as a request for timing-derived state; a positive supplied seed follows a separate path. NetworKit uses a thread-local engine seeded either identically across threads or with the thread identifier added, according to its explicit option. Equal seeds and equal-position draws can produce copied streams, while changing thread count or work scheduling changes how draws are assigned. In a hypothetical two-draw Bernoulli experiment, the variance of their mean is $1/8$ for independent draws but $1/4$ for identical draws. Thread safety alone does not settle this scientific distinction.~\citep{exec25_cnet_snap,exec25_cnet_networkit,exec25_cnet_choice}

Both remain conditional direct-integration candidates for specified graph services. NetworKit's MIT declaration and SNAP's BSD-style conditions support concrete artifact review, while dependency rights and installed behavior remain open. A port should be justified by a required runtime or representation constraint and tested against the same graph receipts. There is no reason to rewrite these graph algorithms merely to claim an internally owned causal library.~\citep{exec25_cnet_networkit,exec25_cnet_snap,exec25_cnet_choice}

\subsection{Repast HPC: distributed copies are not independent borrowers}

Repast HPC represents a shared network through local agents, nonlocal copies and complementary edges across processes. Its agent identity consists of original identifier, starting process and type; the current process is excluded from equality. That supports migration without changing identity, but it also means a bank evaluation must count scientific units independently of process-local copies. The selected graph stores one edge per ordered endpoint pair; a duplicate replaces the edge when allowed or marks the existing edge as conflicted. This behavior is different from a transaction multigraph and needs an explicit receipt.~\citep{exec25_cnet_repast}

The directed-neighbor delegate appends incoming and outgoing neighbors into the output vector. A reciprocal neighbor can appear twice, because the shared item helper does not deduplicate. This differs from the set-based generic adjacency inspected in Repast Simphony. Porting a model between the two frameworks therefore requires a semantic comparison rather than matching method names. Directional, unique-neighbor and incident-edge exposure should be separate adapter operations.~\citep{exec25_cnet_repast,exec25_cnet_choice}

The schedule comparator orders events by tick only, and the execution loop invokes same-tick callbacks sequentially. There is no explicit secondary phase key in that comparator. A repeating event with zero interval can re-enter the same tick indefinitely under the inspected loop, so positive finite intervals and future-time progress should be enforced at the boundary. Across processes, the runner selects the global minimum next tick. That synchronization does not itself refresh agent attributes or network projections. The source provides separate synchronization operations and expressly notes that synchronizing agent state does not update projection information.~\citep{exec25_cnet_repast}

Seed initialization is equally specific. A global-seed property can intentionally give processes the same seed, while the local-seed path derives rank-dependent state from the configured seed and records the effective value. These options serve different coupling purposes. Migration, rank count and assignment of draws to agents must be recorded if a public experiment is expected to replay. Agent ownership, graph-projection freshness and random-state ownership together define the simulated transition law. A valid MPI exchange is not evidence of causal or statistical independence.~\citep{exec25_cnet_repast,exec25_cnet_choice}

The decision is conditional use for distributed known-law experiments when scale justifies MPI, serialization and synchronization costs. Its BSD-style core terms do not clear every Boost, NetCDF or other installed dependency. Qualification should compare one-process and multi-process scientific outputs under a deliberately fixed coupling, exercise migration and stale-projection faults, and record refused or stalled requests. Small graph enumeration should precede expensive distributed runs.~\citep{exec25_cnet_repast,exec25_cnet_choice}

\subsection{igraph's C interface is a distinct interoperability option}

igraph is included explicitly as a C library usable from a C++ system. The selected neighbor API exposes direction, loop treatment and parallel-edge multiplicity, including a choice to collapse multiple edges. These are positive capabilities for a causal exposure adapter. The exact signature belongs to the pinned revision; a wrapper must not assume that an older installed version has the same interface. Vertex deletion reindexes IDs and invalidates iterators, and the induced-subgraph APIs expose mappings. Borrower identity should therefore live outside transient integer positions and be restored through the recorded maps.~\citep{exec25_cnet_igraph}

Its selected Dijkstra implementation validates weights, rejects negative values and NaNs, and documents explicit unreachable distances. That is useful failure behavior, but a caller must still decide whether a monetary relationship weight is a distance, a strength or an inverse cost. The proposed request includes units and transformation rules. A strong exposure relationship is not automatically a long path, and changing this interpretation can reverse which borrowers are considered close.~\citep{exec25_cnet_igraph,exec25_cnet_choice}

Random ownership requires attention even where a variable is thread-local. The selected default-pointer declaration is thread-local, but it initially points to a shared default generator object and backing state. Setting a custom default changes the pointer and requires that object to remain alive. A caller cannot infer private generator ownership merely from the pointer's storage class. The proposed native fixture should install and retain task-owned generators, check concurrent use and restore prior defaults deliberately. This is a source-based ownership finding, not a reported concurrency benchmark.~\citep{exec25_cnet_igraph}

The C interface can be adopted directly where its representation and failure contracts fit. Its GPL 2-or-later source declaration requires a distribution review for the intended product; porting or placing a process boundary around it is not declared an automatic rights solution. Complete intended artifacts, dependencies and native interface behavior remain qualification gates.~\citep{exec25_cnet_igraph,exec25_cnet_choice}

\subsection{A common public experiment contract}

The CLI and MCP surface should bind the scientific target, exposure model, graph and node-map hashes, edge units and multiplicity, event and availability clocks, partition generator, realized partition, conditional or marginal probabilities, random streams, worker layout and numerical tolerances. The response should preserve exposed-arm counts, probability uncertainty, graph transformations, typed failures, unsupported targets and complete cost denominators. An LLM may help propose this specification, but an unreviewed tool choice must not convert a graph service into a causal estimator or silently choose the empty-arm convention.~\citep{exec25_cnet_choice}

The public qualification sequence is concrete. First enumerate tiny graphs and all assignment cases, comparing marginal and joint exposure probabilities, HT moments and ratio failure behavior. Next test representation changes: duplicated transactions, removed nodes, compacted identifiers, weighted restoration and directed paths. Then freeze random-state and scheduling rules across threads or ranks, including stale graph caches and stale distributed projections. Finally compare matched native implementations with all attempts, failures, refusals and accepted outputs recorded. These tasks can proceed independently of unavailable bank truth, while expensive bank trials remain conditional on demonstrated public merit.~\citep{exec25_cnet_choice}

With a hypothetical USD 30,000 study budget, 1,000 attempts, 600 completions and 400 accepted outputs imply costs of USD 30 per attempt, USD 50 per completion and USD 75 per accepted output. None is a measured return or a price quote. They illustrate why a fast executable with frequent invalid exposure results can be economically inferior to a slower method with an appropriate scientific contract. The known repair routes are exact-law testing, explicit graph and random-state custody, typed refusal, rights review and matched native evaluation. Unknown bank interference, unsupported assignment truth and unrestricted dependent-data inference still require information or justified restrictions; this review does not establish complete C++ network qualification.~\citep{exec25_cnet_choice}
